\documentclass[11pt]{article}
\usepackage[margin=2.5cm]{geometry}
\usepackage{amsmath,amssymb,amsthm,longtable,booktabs}
\usepackage[hidelinks]{hyperref}
\newtheorem{theorem}{Theorem}
\newtheorem{lemma}[theorem]{Lemma}
\newtheorem{corollary}[theorem]{Corollary}
\newtheorem{proposition}[theorem]{Proposition}
\theoremstyle{definition}
\newtheorem{definition}[theorem]{Definition}
\newtheorem{remark}[theorem]{Remark}
\newcommand{\Z}{\mathbb{Z}}
\newcommand{\F}{\mathbb{F}}
\newcommand{\R}{\mathbb{R}}
\newcommand{\tr}{\operatorname{tr}}
\newcommand{\Min}{\operatorname{Min}}
\setlength{\emergencystretch}{3em}

\title{No $421$ equiangular lines in $\R^{44}$, no $94$ in $\R^{20}$, no $74$ in $\R^{19}$ and no $59$ in $\R^{18}$}
\author{EQ-UP (upper-bound agent of the equiangular-lines team)\\
\small Status: Theorem~\ref{thm:main} is rated VALID by the auditor EQ-AUD (A.10) and by the referee EQ-REF1.\\
\small Theorem~\ref{thm:d20} is rated VALID by EQ-AUD (A.19) and by the referee EQ-REF4 (text fixes applied).\\
\small Theorem~\ref{thm:d19} is rated VALID by EQ-AUD (A.28, no fix); the referee has not reported yet.\\
\small Theorem~\ref{thm:d18} is submitted to EQ-AUD (24 October 2026, 08:37 UTC) and not yet audited.\\
\small Nothing here is to be reported as a result before both the auditor and a referee rate it VALID.}
\date{24 October 2026}

\begin{document}
\maketitle

\begin{abstract}
Let $N(d)$ be the maximum number of equiangular lines in $\R^d$. We prove that there are no $421$
equiangular lines in $\R^{44}$, no $94$ in $\R^{20}$, no $74$ in $\R^{19}$ and no $59$ in $\R^{18}$, so
\[
344\le N(44)\le 420,\qquad 90\le N(20)\le 93,\qquad 72\le N(19)\le 73,\qquad 57\le N(18)\le58 .
\]
The published upper bounds are $422$ \cite[Table~3]{BY14}, \cite[Thm.~8]{KT19}, $94$ and $74$
\cite[Thms.~5.5, 5.2]{GSY21}, and $59$ \cite{GS24}. All four proofs follow the method of Greaves,
Syatriadi and Yatsyna.

For $d=44$ the $2$-adic structure of Seidel characteristic polynomials forces
$\chi_S(x)=(x+7)^{377}(x-59)^{16}(x-61)^{15}Q(x-59)$ with $Q$ of degree $13$. There are exactly $82$
possibilities for $Q$. Each one is excluded by an exact Farkas certificate showing that the interlacing
conditions on the $421$ vertex-deleted submatrices are infeasible. Two independently written programs
(Python/FLINT and PARI/GP) carried out the computation.

For $d=20$ the reduction of \cite{GSY21} leaves $169$ candidate spectra. Interlacing with the
$2$-adic class condition of \cite{GSY23} leaves $20$. An analysis of pairs of vertices (compatibility
at the simple eigenvalues, Jacobi's identity for the adjugate, and linear systems around one vertex) leaves only
$(x-21)(x-19)^{16}(x-15)^3(x+5)^{74}$. For this spectrum the components of the lines in the
eigenspaces of $21$ and $15$ are rigid. The question becomes a finite problem about vectors in the
$16$-dimensional eigenspace of $19$, which we decide by exhaustive search. Two programs without
shared code exclude every case. One of them produces an exact integer Farkas certificate for every
leaf of its search, and all certificates are re-checked without an LP solver.

For $d=19$ the same method reduces $74$ lines to the spectrum
$(x-17)(x-15)^{15}(x-11)^3(x+5)^{55}$, and both programs exclude every case; again every
certificate of the certificate-producing program is re-checked.

For $d=18$, $59$ lines force $\chi_S=(x+5)^{41}(x-11)\,\psi(x-11)$ with $\psi$ of degree $17$. Two
programs enumerate the $\psi$ that satisfy the $2$-adic conditions of Greaves and Yatsyna, on
prefixes of the coefficient vector, and find the same $1{,}214$ candidates. Interlacing with these
conditions, the pair analysis at odd order, and a third level (triangle types of three vertices,
linear systems for each edge type, and a relaxed vertex-star test for large cliques) exclude all
but one. The last one, $(x+5)^{41}(x-11)^{14}(x-15)^2(x^2-21x+94)$, forces a rigid structure on
the $40$ vertices of one deck member, and a finite computation excludes it. Both programs, written
independently, exclude every candidate. As a by-product we record a parity sharpening of the
relative bound. It comes from known tools.
\end{abstract}

\section{Results}

A set of $n$ equiangular lines in $\R^d$ with common angle $\arccos\alpha$ is given by unit
vectors $v_1,\dots,v_n$ with Gram matrix $G=I+\alpha S$, where $S$ is a \emph{Seidel matrix}:
symmetric, with zero diagonal and off-diagonal entries $\pm1$. Since $G\succeq0$ and
$\operatorname{rank}G=\dim\operatorname{span}\{v_i\}\le d$, the smallest eigenvalue of $S$ is at
least $-1/\alpha$, with multiplicity $n-\operatorname{rank}G$. $N_\alpha(d)$ denotes the maximum
for a fixed $\alpha$, and $\chi_M(x)=\det(xI-M)$.

\begin{theorem}\label{thm:main}
There is no set of $421$ equiangular lines in $\R^{44}$. Consequently $N(44)\le420$.
\end{theorem}

\begin{theorem}\label{thm:d20}
There is no set of $94$ equiangular lines in $\R^{20}$. Consequently $N(20)\le93$.
\end{theorem}

\begin{theorem}\label{thm:d19}
There is no set of $74$ equiangular lines in $\R^{19}$. Consequently $N(19)\le73$.
\end{theorem}

\begin{theorem}\label{thm:d18}
There is no set of $59$ equiangular lines in $\R^{18}$. Consequently $57\le N(18)\le58$.
\end{theorem}

Theorems~\ref{thm:d20}, \ref{thm:d19} and~\ref{thm:d18} are computer-assisted. Theorems~\ref{thm:d20}
and~\ref{thm:d19} share the chain of reductions of \S\S\ref{sec:cand}--\ref{sec:model}, and the final
search of \S\ref{sec:comp35} is done by two independently written programs, one of which writes an
exact certificate for every leaf; Section~\ref{sec:d19} lists the computations for both theorems.
Theorem~\ref{thm:d18} (\S\ref{sec:59}) runs the chain at the odd order $59$, with a third level and a
structural argument for one spectrum; every step is done by two independently written programs.

The bound $N(44)\le421$ is the starting point. It is not stated in the literature, but it follows
directly from \cite[Thm.~5.21 and Cor.~3.6]{GKMS16}, or from \cite[Lemma~3.1]{GSY21} with
Cauchy--Schwarz, once $N(43)=344$ forces full rank. Within our team it is claim C1 of EQ-LIT,
which EQ-AUD rated \emph{valid with fixes} (attribution). For completeness, \S\ref{sec:422} gives a
proof in the notation used here.

\paragraph{Previous bounds.} $N(43)=344$, by Taylor's unitary two-graph \cite{Taylor} and the
relative bound. For $d=44$ the relative bound gives $\lfloor48\cdot44/5\rfloor=422$. The same
value is the best upper bound in \cite[Table~3]{BY14} (SDP bounds) and \cite[Thm.~8]{KT19}; EQ-AUD
found no published bound below $422$. The lower bound $N(44)\ge344$ is only the $43$-dimensional
system. For $d=19$ and $20$ the published ranges are $72\le N(19)\le74$ and $90\le N(20)\le94$
\cite[Thms.~5.2, 5.5]{GSY21}, \cite[Table~1]{GS24}. The lower bounds come from Asche's $72$ lines in
$\R^{19}$ \cite{Szo19} and Taylor's $90$ lines in $\R^{20}$ \cite{Tay71}. For $d=18$ the published
range is $57\le N(18)\le59$: the lower bound is the lattice construction of \cite{GSY23}, and the
upper bound is \cite{GS24}, which excludes $60$ lines. Greaves and Syatriadi
\cite{GS24} note that $N(18)$, $N(19)$ and $N(20)$ are the only unknown values for $d\le41$.

\paragraph{By-product.} Proposition~\ref{prop:U} (\S\ref{sec:U}) combines
\cite[Lemma~3.1]{GSY21} with Cauchy--Schwarz \cite[Thm.~5.6]{GKMS16}. Applied at the top value
$n_0=\lfloor R\rfloor$ of a non-integral relative bound $R$, it gives
$N_{1/a}(d)\le\lfloor R\rfloor-1$ for many $(a,d)$, including the infinite family
$N_{1/a}(a^2-5)\le\lfloor(a^2-1)(a^2-5)/5\rfloor-1$ for $a\equiv\pm3\pmod{10}$, $a\ge7$.
These exclusions are not new in substance. On all $140$ cases we tabulated, they coincide exactly
with what the even-centre inequality of \cite[Thm.~5.21]{GKMS16} gives, and the family also follows
from that theorem. We have not seen the family stated.

\section{Tools}

\subsection{Known results used}
\begin{enumerate}
\item[(K1)] (Neumann, \cite[Thm.~3.4]{LS73}; see also \cite[Thm.~1.2]{BY14}) If $n>2d$ lines are
equiangular in $\R^d$ with angle $\arccos\alpha$, then $1/\alpha$ is an odd integer.
\item[(K2)] (Lemmens--Seidel \cite[Thm.~4.5]{LS73}; see also \cite{KT19}) $N_{1/3}(d)=2(d-1)$ for
$d\ge15$.
\item[(K3)] $N_{1/5}(d)\le276$ for $23\le d\le59$ \cite[Thm.~1.1]{Yu17} (and for $d\le185$,
\cite{CKLY22}).
\item[(K4)] (Relative bound \cite{vLS66,LS73}) If $a^2>d$ then
$N_{1/a}(d)\le(a^2-1)d/(a^2-d)=d+d(d-1)/(a^2-d)$.
\item[(K5)] (Cauchy interlacing) If $S[i]$ is $S$ with row and column $i$ deleted, the eigenvalues of
$S[i]$ interlace those of $S$; and $\sum_{i=1}^n\chi_{S[i]}(x)=\chi_S'(x)$.
\end{enumerate}

\subsection{\texorpdfstring{The $2$-adic structure of $\chi_S$}{The 2-adic structure of chi\_S}}
Write $S=J-I-2A$, where $A$ is the adjacency matrix of a graph on $n$ vertices. The facts below are
due to Greaves and Yatsyna \cite[\S3]{GY19} and Greaves, Syatriadi and Yatsyna
\cite[Thm.~2.5, Lemmas~2.7, 2.8]{GSY21}. We include short proofs because we use a parity refinement
for an arbitrary odd shift (Lemma~\ref{lem:type2}) and its transfer to factors
(Lemma~\ref{lem:divide}(2)).

\begin{definition}[{\cite[Def.~2.6]{GSY21}}]
A monic $P=\sum_i a_ix^{N-i}\in\Z[x]$ of degree $N$ is \emph{type 2} if $2^i\mid a_i$ for all $i$,
equivalently $P(2y)/2^N\in\Z[y]$; it is \emph{weakly type 2} if $2^{i-1}\mid a_i$ for all $i\ge1$,
equivalently $P(2y)/2^{N-1}\in\Z[y]$.
\end{definition}

\begin{lemma}\label{lem:match}
For a graph adjacency matrix $A$ of order $n$, $\det(zI+A)\equiv\sum_{j\ge0}m_jz^{n-2j}\pmod 2$,
where $m_j$ is the number of $j$-edge matchings.
\end{lemma}
\begin{proof}
Modulo $2$ signs disappear, so $\det(zI+A)\equiv\sum_\sigma\prod_i(zI+A)_{i\sigma(i)}$. The terms of
$\sigma$ and $\sigma^{-1}$ coincide because $zI+A$ is symmetric, so all $\sigma\ne\sigma^{-1}$ cancel in
pairs. An involution with $j$ transpositions contributes $z^{n-2j}\prod A_{uv}$ over its
transpositions $(u\,v)$. This is $z^{n-2j}$ if they form a matching of the graph and $0$ otherwise.
\end{proof}

\begin{lemma}[{cf.\ \cite[Thm.~2.5, Lemma~2.7]{GSY21}}]\label{lem:type2}
Let $S$ be a Seidel matrix of order $n$, $\kappa$ odd, $P(x)=\chi_{S-\kappa I}(x)=\chi_S(x+\kappa)$ and
$\widetilde P(y)=P(2y)/2^{n-1}$. Then $\widetilde P\in\Z[y]$. If $n$ is even, $\widetilde P\equiv0\pmod 2$,
i.e.\ $P$ is type 2. If $n$ is odd, $\widetilde P\bmod 2$ is a polynomial in $y^2$.
\end{lemma}
\begin{proof}
Put $k=(\kappa+1)/2\in\Z$ and $B=A+kI$, so that $S-\kappa I=J-2B$. By the matrix determinant lemma,
\[
P(2y)=\det\bigl(2(yI+B)-\mathbf 1\mathbf 1^{\mathsf T}\bigr)=2^nD(y)-2^{n-1}\,\mathbf 1^{\mathsf T}\operatorname{adj}(yI+B)\mathbf 1,
\qquad D(y)=\det(yI+B).
\]
Hence $\widetilde P=2D-\mathbf 1^{\mathsf T}\operatorname{adj}(yI+B)\mathbf 1\in\Z[y]$. Since $\operatorname{adj}(yI+B)$ is
symmetric, the off-diagonal entries contribute an even amount, so
$\widetilde P\equiv\tr\operatorname{adj}(yI+B)=D'(y)\pmod 2$. By Lemma~\ref{lem:match},
$D(y)=\det((y+k)I+A)\equiv\sum_jm_j(y+k)^{n-2j}$. If $n$ is even, all exponents are even, so
$D'\equiv0$. If $n$ is odd, $D'(y)\equiv\sum_jm_j(y+k)^{n-2j-1}$, a polynomial in
$(y+k)^2\equiv y^2+k^2$.
\end{proof}

\begin{lemma}[{cf.\ \cite[Lemma~2.8]{GSY21}}]\label{lem:divide}
Let $P=FG$ with $P,F,G\in\Z[x]$ monic of degrees $N,f,g$.
\begin{enumerate}
\item A monic $P$ is type 2 iff $\rho/2$ is an algebraic integer for every root $\rho$. Hence every
monic integer factor of a type-2 polynomial is type 2.
\item If $F$ is type 2 and $P$ is weakly type 2, then $G$ is weakly type 2. If moreover
$\widetilde P:=P(2y)/2^{N-1}\bmod2$ and $\widehat F:=F(2y)/2^f\bmod 2$ are polynomials in $y^2$, then so
is $\widetilde G:=G(2y)/2^{g-1}\bmod 2$.
\end{enumerate}
\end{lemma}
\begin{proof}
(1) $P(2y)/2^N=\prod(y-\rho/2)$, and its coefficients are symmetric functions of the $\rho/2$.
(2) $\widetilde P=\widehat F\,\widetilde G$ with $\widehat F\in\Z[y]$ monic and $\widetilde P\in\Z[y]$, so long
division by a monic integer polynomial gives $\widetilde G\in\Z[y]$. Write
$\widetilde G\equiv G_0(y^2)+yG_1(y^2)$ and $\widehat F\equiv H(y^2)$ in $\F_2[y]$, with $H\neq0$ because
$\widehat F$ is monic. Then $\widetilde P\equiv HG_0(y^2)+yHG_1(y^2)$ is a polynomial in $y^2$ only if
$HG_1=0$, i.e.\ $G_1=0$.
\end{proof}

\begin{remark}
The hypothesis on $\widehat F$ in Lemma~\ref{lem:divide}(2) cannot be dropped. For
$P=(x-2)(x^2-x+2)=x^3-3x^2+4x-4$ one has $\widetilde P=2y^3-3y^2+2y-1\equiv y^2+1$, but for
$G=x^2-x+2$ one has $\widetilde G=2y^2-y+1\equiv y+1\pmod2$; here $\widehat F=y-1$.
\end{remark}

\begin{lemma}[forcing; the step in the proof of {\cite[Lemma~3.1]{GSY21}}]\label{lem:force}
Let $S$ be a Seidel matrix of order $n$, $\gamma=n\bmod2$, and $\kappa$ odd. Split the eigenvalues
(with multiplicity) as $\Lambda=\Lambda_0\sqcup E$, where every element of $\Lambda_0$ is an integer
$\equiv\kappa\pmod 2$. Let $\eta=|E|\ge1$ and $T=\sum_{\lambda\in E}(\lambda-\kappa)^2$. If
$\kappa\notin E$, then $\eta^\eta4^{\eta-\gamma}\le T^\eta$. In particular, if
$T^\eta<\eta^\eta4^{\eta-\gamma}$ then $\kappa\in E$.
\end{lemma}
\begin{proof}
$\chi_S(x+\kappa)=F(x)q(x)$ with $F=\prod_{\lambda\in\Lambda_0}(x-(\lambda-\kappa))$, which is type 2
because its roots are even integers, and $q=\prod_{\lambda\in E}(x-(\lambda-\kappa))$. By
Lemmas~\ref{lem:type2} and~\ref{lem:divide}, $q$ is weakly type 2, and type 2 if $n$ is even. So its
constant term $\pm\prod_E(\lambda-\kappa)\ne0$ is divisible by $2^{\eta-\gamma}$. By AM--GM,
$(T/\eta)^\eta\ge\prod_E(\lambda-\kappa)^2\ge4^{\eta-\gamma}$.
\end{proof}

\section{\texorpdfstring{No $422$ lines (known tools)}{No 422 lines (known tools)}}\label{sec:422}

\begin{proposition}[{EQ-LIT claim C1; \cite[Thm.~5.21, Cor.~3.6]{GKMS16}, \cite[Lemma~3.1]{GSY21}}]\label{prop:422}
There is no set of $422$ equiangular lines in $\R^{44}$.
\end{proposition}
\begin{proof}
\emph{Reduction.} Suppose there are $n\ge422$ lines in $\R^{44}$; keep $422$ of them. By (K1)
($422>88$) the angle is $\arccos\frac1a$ with $a$ odd. By (K2) ($2\cdot43=86$), by (K3), and by (K4)
for $a\ge9$ ($R(a,44)\le R(9,44)=3520/37<96$), we get $a=7$. If the lines spanned at most $43$
dimensions, (K4) would give $422\le R(7,43)=344$. So $S$ (order $422$) has smallest eigenvalue $-7$ of
multiplicity exactly $378$, and its other eigenvalues $\lambda_1,\dots,\lambda_{44}$ satisfy
$\sum\lambda_i=7\cdot378=2646$ and $\sum\lambda_i^2=422\cdot421-49\cdot378=159140$ (from $\tr S=0$,
$\tr S^2=n(n-1)$).

\emph{Exclusion.} Put $\kappa=61$, $x_i=\lambda_i-61$: $\sum x_i=2646-44\cdot61=-38$ and
$\sum x_i^2=159140-122\cdot2646+44\cdot61^2=52$. Let $E$ be the set of $i$ with $x_i\ne0$ and $\eta=|E|$.
Lemma~\ref{lem:force} with $\Lambda_0$ the $378$ eigenvalues $-7$ and the eigenvalues equal to $61$
(all odd), and $\gamma=0$, gives $4\eta\le52$, so $\eta\le13$. By Cauchy--Schwarz,
$38^2\le\eta\cdot52\le676$, which is false.
\end{proof}

The same reduction applies to every $n$ with $344<n\le422$: the angle is $\arccos\frac17$ and $-7$
has multiplicity exactly $n-44$.

\section{Proof of Theorem~\ref{thm:main}}\label{sec:421}

Suppose there are $421$ equiangular lines in $\R^{44}$. As above, $S$ has eigenvalue $-7$ of
multiplicity exactly $377$, and the other $44$ eigenvalues satisfy $\sum\lambda=2639$ and
$\sum\lambda^2=158347$.

\paragraph{Step 1: stripping.} Apply Lemma~\ref{lem:force} repeatedly. At each step $\Lambda_0$
consists of the $377$ eigenvalues $-7$ and the eigenvalues already identified, and $E$ of the $r$
remaining ones; $\sum_E\lambda$ and $\sum_E\lambda^2$ are known at every stage. Sixteen steps with
$\kappa=59$, followed by fifteen with $\kappa=61$, each satisfy $T^r<r^r4^{r-1}$ ($\gamma=1$). The
first step has $r=44$, $T=\sum(\lambda-59)^2=109$, and $109^{44}<44^{44}4^{43}$. Removing an
eigenvalue equal to $\kappa$ leaves $T$ unchanged; after the $\kappa=59$ steps, $r=28$ and
$\sum_E(\lambda-61)^2=49$; the last step is $r=14$, where $49^{14}<14^{14}4^{13}$, while at $r=13$
the test fails. All $31$ inequalities are checked in both programs. Hence
\[
\chi_S(x)=(x+7)^{377}(x-59)^{16}(x-61)^{15}\,Q(x-59),
\]
where $Q$ is monic of degree $13$ and real-rooted, and its roots $x_j=\lambda_j-59$ satisfy
$\sum x_j=13$ and $\sum x_j^2=49$. So $Q=x^{13}-13x^{12}+60x^{11}+\cdots$.

\paragraph{Step 2: the lattice for $Q$.} Apply Lemmas~\ref{lem:type2} and~\ref{lem:divide} with
$\kappa=59$. Here $P=\chi_S(x+59)=F\,Q$ with $F=(x+66)^{377}x^{16}(x-2)^{15}$, which is type 2, and
\[
\widehat F=(y+33)^{377}y^{16}(y-1)^{15}\equiv(y+1)^{392}y^{16}=\bigl((y+1)^{196}y^{8}\bigr)^2\pmod 2,
\]
a polynomial in $y^2$. Since $n=421$ is odd, $\widetilde P\bmod2$ is a polynomial in $y^2$. Hence
$\widetilde Q=Q(2y)/2^{12}=\sum_i a_i2^{1-i}y^{13-i}\in\Z[y]$, and its coefficients at odd powers of $y$
are even. In terms of $Q=\sum_ia_ix^{13-i}$:
\[
2^{i-1}\mid a_i\ (i\ \text{odd}),\qquad 2^{i}\mid a_i\ (i\ \text{even}).
\]
The first condition is the weakly-type-2 property \cite[Lemma~2.7]{GSY21}. The second comes from
\cite[Lemma~3.8]{GY19} ($2^r\mid a_r$ for even $r$; for odd order this is the parity refinement of
Lemma~\ref{lem:type2}), transferred to $Q$ by Lemma~\ref{lem:divide}(2). The proof does not depend
on the second condition; see the paragraph on robustness in \S\ref{sec:comp}.

\paragraph{Step 3: enumeration.} There are exactly $82$ monic real-rooted $Q\in\Z[x]$ of degree $13$
with $a_1=-13$, $a_2=60$ and these divisibilities (\S\ref{sec:comp}). Table~\ref{tab:cand}
(Appendix~\ref{app:cand}) lists their factorisations. For each candidate $p=\chi_S$, the number
$\delta$ of distinct eigenvalues satisfies $5\le\delta\le8$.

\paragraph{Step 4: interlacing characteristic polynomials.} This is the method of
\cite[\S5]{GSY23}; the term \emph{deck} is from \cite[Def.~2.5]{GS24}, and \cite{GS25} treats the
enumeration of interlacing polynomials by linear programming. Fix a candidate $p=\chi_S$. Let $\Min_p$ be the product of the distinct monic
irreducible factors of $p$, $\delta=\deg\Min_p$, $\mu=p/\Min_p$ and $h=p'/\mu$. For each vertex $i$,
(K5) gives $\mu\mid\chi_{S[i]}$: a root of multiplicity $e$ of $p$ has multiplicity at least $e-1$ in
$\chi_{S[i]}$, and the same holds for all its conjugates. So $\chi_{S[i]}=\mu\,g_i$, where $g_i\in\Z[x]$ is
monic of degree $\delta-1$ and lies in the \emph{deck} of $p$, the set of all $g$ with:
\begin{enumerate}
\item[(D1)] $g=x^{\delta-1}+b_1x^{\delta-2}+b_2x^{\delta-3}+\cdots$ with $b_1=c_1$ and $b_2=c_2+n-1$, where
$\Min_p=x^\delta+c_1x^{\delta-1}+c_2x^{\delta-2}+\cdots$. This follows from $\tr S[i]=0$ and
$\tr S[i]^2=(n-1)(n-2)$.
\item[(D2)] $g(\theta_j)/\Min_p'(\theta_j)\ge0$ for every root $\theta_j$ of $\Min_p$, i.e.\ $g$
interlaces $\Min_p$ ($g$ has one root in each $[\theta_k,\theta_{k+1}]$). Indeed
$(xI-S)^{-1}_{ii}=\chi_{S[i]}/\chi_S=g/\Min_p=\sum_jw_j/(x-\theta_j)$ with
$w_j=\|E_je_i\|^2\ge0$, where $E_j$ is the orthogonal projection onto the $\theta_j$-eigenspace, and
the residue at $\theta_j$ is $w_j=g(\theta_j)/\Min_p'(\theta_j)$.
\item[(D3)] $g(x-1)$ is type 2. Indeed $\chi_{S[i]}(x-1)=\chi_{S[i]+I}(x)$ is type 2 by
Lemma~\ref{lem:type2} ($n-1=420$ is even, $\kappa=-1$), and $g(x-1)$ is a monic integer factor of it
(Lemma~\ref{lem:divide}(1)).
\item[(D6)] $g(-7)\ne0$. By interlacing, $S[i]$ has $-7$ with multiplicity at least $376$. It has
multiplicity at most $376$, since otherwise the $420$ remaining lines would span at most $43$
dimensions, contradicting $N_{1/7}(43)\le344$. Since $(x+7)^{376}$ divides $\mu$, $g(-7)\ne0$.
\item[(D7)] Every root of $g$ lies in $[2166/44,\,3098/44]$. By (D6) the roots of $g$ are among the
$44$ eigenvalues of $S[i]$ other than $-7$. Their sum is $7\cdot376=2632$ and their sum of squares is
$420\cdot419-49\cdot376=157556$, so their variance sum is $5040/44$. By Cauchy--Schwarz each of
them lies within $\sqrt{43\cdot5040}/44<466/44$ of the mean $2632/44$.
\end{enumerate}
Moreover, $\sum_i\chi_{S[i]}=p'$ (K5) gives $\sum_{i=1}^{421}g_i=h$, where
$h=\sum_k e_kq_k'\Min_p/q_k$ if $p=\prod_kq_k^{e_k}$ with distinct irreducible $q_k$. Conditions
(D1)--(D3) are those of \cite[\S5]{GSY23}; condition (iv) there, on $\mathcal P_{n-1,7}$, is only
used when $n-1$ is odd. Conditions (D6) and (D7) are specific to this case.

\paragraph{Step 5: Farkas certificates.} For each of the $82$ candidates the deck was computed
exhaustively (at most $73$ members), and an integer vector $Y$ was found with
$\operatorname{Coeff}(g)\cdot Y\ge0$ for every deck member $g$ and $\operatorname{Coeff}(h)\cdot Y<0$,
where $\operatorname{Coeff}$ is the coefficient vector (a certificate of infeasibility in the sense of
\cite[\S5]{GSY23}). Then $0\le\sum_i\operatorname{Coeff}(g_i)\cdot Y=\operatorname{Coeff}(h)\cdot Y<0$, a
contradiction. So no candidate is the characteristic polynomial of a Seidel matrix, and there are no
$421$ equiangular lines in $\R^{44}$. Since every set of more than $421$ lines contains $421$ of them,
$N(44)\le420$. \qed

\section{Computation}\label{sec:comp}

\paragraph{Enumeration principle} (McKee--Smyth, as in \cite[\S2.3]{GSY21}). If
$f=\sum_{j=0}^tc_jx^{t-j}$ is real-rooted, so is every derivative $D_i=f^{(t-i)}$ (of degree $i$),
and $D_i'=D_{i-1}$. Given $c_0,\dots,c_{i-1}$, write $D_i=P_i+c_i(t-i)!$. If
$\mu_1\le\dots\le\mu_{i-1}$ are the roots of $D_{i-1}$, then $D_i$ is real-rooted only if
$(-1)^{i-k}D_i(\mu_k)\ge0$ for all $k$, which confines $c_i$ to an interval. All $c_i$ in that interval
and in the required residue class are tried, and each $D_i$ is then tested exactly. For decks, the
$k$-th root of $D_i$ must in addition lie in $[\theta_k,\theta_{k+t-i+1}]\cap[2166/44,3098/44]$, which
follows from (D2), (D7) and Rolle's theorem.

\paragraph{Program A} (Python 3 with \texttt{python-flint} 0.9 and SciPy/HiGHS). The code is in
\texttt{src/n421\_enum.py}, \texttt{src/n421\_decks.py} and \texttt{src/eqlines/}. Interval endpoints use
ball arithmetic.
Real-rootedness and the root windows are decided by FLINT's certified real root isolation, and signs
at algebraic eigenvalues by exact divisibility or refined balls. Certificates are found by HiGHS and
made exact on the tight set. One candidate (no.~72) needed an exact search over vertices of the dual
polyhedron (\texttt{find\_farkas\_vertex}). Every certificate is verified in exact integer arithmetic.

\paragraph{Program B} (PARI/GP 2.15, \texttt{verify/n421.gp}), written separately and sharing no
code with Program A. From Program A it reads only the stripping chain and the certificates $Y$
(\texttt{data/n421/export.gp}). It re-checks the $31$ stripping inequalities and re-enumerates $Q$
with exact Sturm sequences and exact root counting with multiplicity. It rebuilds each $p$, $\Min_p$
and $h$ from its own factorisation, re-enumerates each deck with exact window and box tests, and
decides (D2) by exact sign evaluation on isolating rational intervals. Finally it checks each
certificate against its \emph{own} deck. Floating point is used only to locate candidate ranges for
coefficients, which are then widened; every accept/reject decision is exact.

\paragraph{Results.} Both programs find the same $82$ polynomials $Q$ ($13{,}784$ search nodes) and
decks of identical sizes ($120{,}527$ deck search nodes in Program B). All $82$ certificates are valid
in both. The largest certificate entry has $35$ decimal digits. Run times: Program A about $20$\,s on
$6$ cores, Program B $77$\,s on one core.

\paragraph{Robustness.} The parity refinement of Step~2 ($2^i\mid a_i$ for even $i$) is not needed.
Requiring only $2^{i-1}\mid a_i$ for all $i\ge1$ (that is, $Q$ weakly type 2), both programs find the
same $190$ polynomials $Q$ ($53{,}127$ search nodes) and decks of identical sizes ($615{,}463$ deck search
nodes, at most $80$ members). Program A finds a certificate for each of the $190$ candidates, and
Program B confirms every one on its own decks (\texttt{--weak} options of Program A,
\texttt{verify/n421\_weak.gp}; Program B takes $395$\,s on one core).

\paragraph{Negative tests.} The same pipeline (\texttt{src/remainder\_case.py}, conditions (D1)--(D7)
with the parameters of each case) was run on two line systems that exist. For $343$ lines in
$\R^{43}$ (the $344$ unitary lines minus one) it finds $704$ remainder polynomials. Nine are not
excluded, and among them is the true spectrum $\{(-7)^{300},49^{42},42\}$. For $125$ lines in
$\R^{21}$ (the $126$ unitary lines for $q=5$ minus one) it finds $14$ remainders. Two are not
excluded, including the true spectrum $\{(-5)^{104},25^{20},20\}$. The moment test of
Proposition~\ref{prop:U} was run on $15{,}565$ parameter pairs coming from all subsets of size
$>N(d-1)$ of the systems with $28$, $126$, $176$, $276$ and $344$ lines, and on actual Seidel matrices
obtained by deleting vertices from the $126$- and $344$-line systems. No existing system was excluded
(\texttt{verify/negative\_tests.py}).

\section{By-product: a parity sharpening of the relative bound}\label{sec:U}

\begin{proposition}[{\cite[Lemma~3.1]{GSY21} with \cite[Thm.~5.6]{GKMS16}}]\label{prop:U}
Let $a\ge3$ be odd and let $S$ be a Seidel matrix of order $n$ whose smallest eigenvalue $-a$ has
multiplicity exactly $n-d$, with other eigenvalues $\lambda_1,\dots,\lambda_d$. For an odd integer
$\kappa$ put
\[
s=\sum_{i=1}^d(\lambda_i-\kappa)=a(n-d)-d\kappa,\qquad
T=\sum_{i=1}^d(\lambda_i-\kappa)^2=n(n-1)-a^2(n-d)-2\kappa a(n-d)+d\kappa^2,
\]
$\gamma=n\bmod2$ and
$\eta^*(T)=\max\{\eta\in\{0,\dots,d\}:\ \eta=0\ \text{or}\ \eta^\eta4^{\eta-\gamma}\le T^\eta\}$. Then
$s^2\le\eta^*(T)\,T$.
\end{proposition}
\begin{proof}
Let $E=\{i:\lambda_i\ne\kappa\}$ and $\eta=|E|$. The remaining eigenvalues are $-a$ (multiplicity $n-d$)
and copies of $\kappa$, all odd integers. If $\eta=0$ then $s=0$. Otherwise Lemma~\ref{lem:force} gives
$\eta\le\eta^*(T)$, and Cauchy--Schwarz gives
$s^2=(\sum_{i\in E}(\lambda_i-\kappa))^2\le\eta\sum_{i\in E}(\lambda_i-\kappa)^2=\eta T$.
\end{proof}

\begin{remark}
For $\gamma=0$, $\eta^*(T)=\min(d,\lfloor T/4\rfloor)$. For $\gamma=1$,
$\eta\,4^{(\eta-1)/\eta}\ge4\eta-4\ln4$ gives $\eta^*(T)\le T/4+\ln4$. Let $\bar\mu=a(n-d)/d$, let
$\delta=\operatorname{dist}(\bar\mu,2\Z+1)$ with $\kappa$ the nearest odd integer, and let
$V=\sum(\lambda_i-\bar\mu)^2=n(a^2-d)(R-n)/d$ with $R=(a^2-1)d/(a^2-d)$. Then $|s|=d\delta$ and
$T=V+d\delta^2$, so for even $n$ the criterion reads (up to the floor) $V<d\delta(2-\delta)$. The
even-centre inequality of \cite[Thm.~5.21]{GKMS16} at the even integer $\mu$ nearest to $\bar\mu$,
with $m=\dim\ker(S-\mu I)=0$ for even $n$ \cite[Cor.~3.6]{GKMS16} and $m\le1$ for odd $n$
($S-\mu I\equiv J-I$ has $\F_2$-rank $n-1$), reads $V+d(1-\delta)^2<d-m$, i.e.\ $V<d\delta(2-\delta)-m$.
For even $n$ the two tests agree up to the floor in $\eta^*$; for odd $n$ they differ by lower-order
terms. Since $V\approx(a^2-1)(R-n)$ and $d<a^2$, either test can only remove the top value
$\lfloor R\rfloor$.
\end{remark}

\begin{corollary}\label{cor:family}
For every odd $a\ge7$ with $a\equiv\pm3\pmod{10}$,
$N_{1/a}(a^2-5)\le\left\lfloor\frac{(a^2-1)(a^2-5)}{5}\right\rfloor-1$.
\end{corollary}
\begin{proof}
Let $d=a^2-5$. As $a^2\equiv4\pmod5$, $(a^2-1)(a^2-5)\equiv2\pmod5$, so $R=n_0+\frac25$ with
$n_0\in\Z$. By (K4), $N_{1/a}(d)\le n_0$. Also $R(a,d-1)=(a^2-1)(a^2-6)/6=R-a^2(a^2-1)/30<n_0-1$, so it
suffices to exclude $n_0$ lines spanning $\R^d$, i.e.\ a Seidel matrix of order $n_0$ with $-a$ of
multiplicity exactly $n_0-d$. Here $V=n_0(a^2-d)(R-n_0)/d=2n_0/d<\frac25(a^2-1)$ and
$\bar\mu=a(a^2-6)/5-2a/(5(a^2-5))$. If $a\equiv3\pmod{10}$ then $a(a^2-6)\equiv9\pmod{10}$, so
$\delta=\frac45-\frac{2a}{5(a^2-5)}\in[0.768,0.8)$ for $a\ge13$. If $a\equiv7\pmod{10}$ then
$a(a^2-6)\equiv1\pmod{10}$ and $\delta=\frac45+\frac{2a}{5(a^2-5)}\in(0.8,0.864]$ for $a\ge7$. Then
$T=V+d\delta^2<0.4a^2+0.7465a^2<1.1465a^2$ and $\eta^*\le T/4+\ln4<0.2867a^2+1.39$, so
\[
\eta^*T<1.147a^2(0.2867a^2+1.39)<0.329a^4+1.6a^2<0.5898(a^2-5)^2\le d^2\delta^2=s^2
\]
for $a\ge7$, and Proposition~\ref{prop:U} excludes $n_0$.
\end{proof}

The family also follows from \cite[Thm.~5.21]{GKMS16} at the even centre (checked exactly for
$7\le a\le2001$ by \texttt{verify/compare\_gkms.py}). Table~\ref{tab:U} (Appendix~\ref{app:U}) lists
all $(a,d)$ with $a\le13$ odd, $d<a^2-2$, non-integral $R$ and $\lfloor R\rfloor>2d$ for which
Proposition~\ref{prop:U} excludes $n_0=\lfloor R\rfloor$: $72$ of the $140$ such pairs. In each row also
$\lfloor R(a,d-1)\rfloor\le n_0-1$, so $N_{1/a}(d)\le n_0-1$. The even-centre test of
\cite[Thm.~5.21]{GKMS16} excludes exactly the same $72$. Only the row $d=44$ ($a=7$) changes a bound
on $N(d)$ itself; in every other row the known bound for another angle is larger.

\section{\texorpdfstring{$94$ lines in $\R^{20}$ and $74$ lines in $\R^{19}$: the candidate spectra}{94 lines in R\textasciicircum 20 and 74 lines in R\textasciicircum 19: the candidate spectra}}\label{sec:cand}

Throughout \S\S\ref{sec:cand}--\ref{sec:d19}, $(n,d,\kappa)$ is $(94,20,19)$ or $(74,19,15)$, and $S$ is
a Seidel matrix of order $n$ that comes from $n$ equiangular lines in $\R^d$.

\begin{lemma}\label{lem:169}
The angle is $\arccos\frac15$, and
\[
\chi_S(x)=(x+5)^{n-d}(x-\kappa)^{d-13}\,2^{13}\psi\Bigl(\frac{x-\kappa}2\Bigr),
\]
where $\psi\in\Z[y]$ is monic and real-rooted of degree $13$ with $\psi=y^{13}+5y^{12}+6y^{11}+\cdots$.
There are exactly $169$ such polynomials $\psi$.
\end{lemma}
\begin{proof}
Since $n>2d$, (K1) gives the angle $\arccos\frac1a$ with $a$ odd. For $a=3$, (K2) allows at most
$2(d-1)<n$ lines. For $a\ge7$, (K4) allows at most $R(7,d)\le R(7,20)=960/29<34$. So $a=5$. If the
lines spanned at most $d-1$ dimensions, (K4) would give $n\le R(5,d-1)$; but $R(5,19)=76<94$ and
$R(5,18)=432/7<74$. So $-5$ has multiplicity exactly $n-d$. For the other $d$ eigenvalues,
$\tr S=0$ and $\tr S^2=n(n-1)$ give
$\sum(\lambda-\kappa)=5(n-d)-d\kappa=-10$ and
$T=\sum(\lambda-\kappa)^2=n(n-1)-25(n-d)-10\kappa(n-d)+d\kappa^2=52$ in both cases.
Lemma~\ref{lem:force} ($\gamma=0$, $\Lambda_0$ the eigenvalues $-5$ and those already found equal to
$\kappa$) applies while $52<4\eta$, i.e.\ for $\eta=d,d-1,\dots,14$. This gives $d-13$ eigenvalues
$\kappa$ and leaves $13$, say $\kappa+2y_j$. Then $\sum y_j=-5$ and $\sum y_j^2=13$, which is the
stated form of $\psi=\prod(y-y_j)$. By Lemma~\ref{lem:type2} ($n$ even), $\chi_S(x+\kappa)$ is
type~2. So is its monic factor $2^{13}\psi(x/2)$ (Lemma~\ref{lem:divide}(1)), i.e.\ $\psi\in\Z[y]$.
The enumeration principle of \S\ref{sec:comp} gives exactly $169$ polynomials $\psi$
(\texttt{data/psi\_deg13\_p1m5\_p2\_13.json}). EQ-AUD (A.11, A.15) and EQ-LIT (claim C12) obtained
the same list with their own programs.
\end{proof}

We number the candidates $0,\dots,168$ as in that file. The two that matter below are
\[
\#3:\ \psi=y^{11}(y+2)(y+3),\qquad \#35:\ \psi=y^9(y-1)(y+2)^3,
\]
that is, $\chi_S=(x+5)^{74}(x-19)^{18}(x-15)(x-13)$ resp.\ $(x+5)^{55}(x-15)^{17}(x-11)(x-9)$ for
\#3, and
\[
\#35:\quad \chi_S=(x-21)(x-19)^{16}(x-15)^3(x+5)^{74}\quad\text{resp.}\quad
(x-17)(x-15)^{15}(x-11)^3(x+5)^{55}.
\]

\paragraph{Level 1: one deleted vertex.} For a candidate $p$ let $\Min_p$, $\delta$, $\mu=p/\Min_p$ and
$h=p'/\mu$ be as in Step~4 of \S\ref{sec:421}, so that $\chi_{S[i]}=\mu g_i$ and $\sum_ig_i=h$. The
order $n-1$ is odd. The deck consists of all monic $g$ of degree $\delta-1$ with:
(D1) and (D2) as in \S\ref{sec:421};
(D3) $g(x-1)$ is weakly type~2, and type~2 if $\mu(x-1)$ is not \cite[Lemmas~2.7, 2.8]{GSY21};
(D4) $\chi_{S[i]}(x-1)=\mu(x-1)g(x-1)$ satisfies \cite[Thm.~2.5]{GSY21};
(D6) $g(-5)\ne0$, since otherwise $S[i]$ would have $-5$ with multiplicity $n-d$ and $n-1$ lines would
lie in $\R^{d-1}$, which is impossible as above;
(iv) $\mu g\bmod 2^7$ is the residue of the characteristic polynomial of some Seidel matrix of order
$n-1$ \cite[\S5, condition~(iv)]{GSY23}.
For (iv) we use the complete set $\mathcal P_{n-1,7}$ of such residues. It has $2048$ elements for both
orders. It is complete because the necessary conditions of \cite[\S3]{GY19} (Remark~3.3 and
Lemmas~3.1, 3.8, 3.12 for $k\le3$) allow exactly $2048$ residues, and for each of them we have an
explicit witness matrix. For each candidate the LP $\sum_gn_g\,\mathrm{Coeff}(g)=\mathrm{Coeff}(h)$,
$n\ge0$, over the deck is decided, with an exact integer Farkas certificate when it is infeasible.
This excludes $147$ of the $169$ candidates for $74$ lines and $149$ for $94$ lines. The survivors
are $\#3,9,20,27,35,42,50,78,81,83,85,89,91,99,100,110,111,142,144,149,153,162$ for $74$ lines,
and the same without $\#42$ and $\#85$ for $94$ lines. (EQ-AUD A.15 reproduced every deck,
verdict and certificate by two methods of its own, and EQ-LIT's \texttt{deck\_dp.py} is a third
program.)

\paragraph{Level 2: pairs of vertices.} Three exact steps, each run by our programs and by EQ-LIT's
(claims C12 and C14; EQ-AUD A.16 reproduced all three).
\begin{enumerate}
\item \emph{Lemma~3.12 of \cite{GY19} and compatibility.} Each deck member must also satisfy the
congruences of \cite[Lemma~3.12]{GY19} at the odd order $n-1$. Let $\mathrm{Sim}$ be the product of
$x-\lambda$ over the simple eigenvalues, $M=\Min_p/\mathrm{Sim}$ and $s=\deg\mathrm{Sim}$. Let $i\ne j$
be vertices with members $g=g_i$ and $g'=g_j$ (possibly $g=g'$). At a simple eigenvalue $\lambda$,
$E_\lambda=u_\lambda u_\lambda^{\mathsf T}$ with $u_\lambda(i)^2=g(\lambda)/\Min_p'(\lambda)$
\cite[Lemma~4.4]{GSY23} and $\Min_p'(\lambda)=\mathrm{Sim}'(\lambda)M(\lambda)$. Since $x^kM$ vanishes
at the multiple eigenvalues, $S^kM(S)=\sum_{\lambda\ \mathrm{simple}}\lambda^kM(\lambda)E_\lambda$.
Hence, for some signs $d_\lambda$ independent of $k$,
\[
N_k:=\sum_{\lambda\ \mathrm{simple}}d_\lambda\lambda^k\,
\frac{\sqrt{g(\lambda)g'(\lambda)}}{\mathrm{Sim}'(\lambda)}=\bigl(S^kM(S)\bigr)_{ij}\in\Z
\qquad(0\le k<s)
\]
(for $k=0$ this is \cite[Lemma~4.7]{GSY23} with $\xi=\mathrm{Sim}$), and $N_k\equiv(x^kM)'(1)\pmod2$.
The parity condition is our addition: \cite{GSY23} impose one only for odd order. Proof: $n$ is even, so
$S\equiv I+J$ and $J^2=nJ\equiv0\pmod2$. Hence $S^k\equiv I+kJ$, and every off-diagonal entry of $F(S)$
is $\equiv F'(1)$ for $F\in\Z[x]$. We call $g,g'$ \emph{compatible} if some signs $d_\lambda$ satisfy
both conditions, and say that $g$ has a \emph{loop} if $g,g$ are compatible. The members used by $S$
form a clique of the compatibility graph, and a member without a loop is used at most once. The LP of
level~1 restricted to each maximal clique leaves $\#3,20,35,50,89,91,99,111,142,153$ (and also
$\#162$ for $94$ lines).
\item \emph{Jacobi pair polynomials.} For $i\ne j$ put $\varphi=\operatorname{adj}(xI-S)_{ij}$. Then
$\varphi=\mu\,\psi_{ij}$ with $\psi_{ij}\in\Z[x]$ of degree $\delta-2$ and leading coefficient $S_{ij}$,
and Jacobi's identity
$\chi_{S[i]}\chi_{S[j]}-\chi_S\chi_{S[ij]}=\varphi^2$ (used in this way in \cite{GS24}, after de Caen)
gives $\chi_{S[ij]}=\mu_2P_{ij}$ with $P_{ij}=(g_ig_j-\psi_{ij}^2)/\mathrm{Sim}$ and $\mu_2=\mu/M$.
We impose: (J1) the coefficients $(S^k)_{ij}$ of $\varphi/p=\sum_k(S^k)_{ij}x^{-k-1}$ satisfy
$(S^k)_{ij}\equiv k\pmod2$, by $S^k\equiv I+kJ$; (J2) $\mathrm{Sim}\mid g_ig_j-\psi_{ij}^2$; (J3)
$\psi_{ij}(\lambda)^2\le g_i(\lambda)g_j(\lambda)$ at the multiple eigenvalues (Cauchy--Schwarz for
$E_\lambda\succeq0$); (J4) $\chi_{S[ij]}$ is real-rooted and interlaces $\chi_{S[i]}$ and
$\chi_{S[j]}$; (J5) $\chi_{S[ij]}$ has the trace, the trace of the square and the $2$-adic form
\cite[Thm.~2.5]{GSY21} of a Seidel matrix of order $n-2$. Every $\psi_{ij}$ allowed by (J1)--(J5) is
enumerated. The graph ``some pair polynomial exists'' and its cliques exclude $\#20$.
\item \emph{Vertex stars.} Since $\sum_{j\ne i}\chi_{S[ij]}=\chi_{S[i]}'$, for every vertex $i$ we get
$\sum_{j\ne i}P_{ij}=(\mu g_i)'/\mu_2$. With unknowns $n_g$ (the number of vertices with member $g$) and
$m^f_{g,P}$ (for a vertex with member $f$, the number of other vertices $j$ with member $g$ and pair
polynomial $P$), this gives the linear system $\sum_gn_gg=h$,
$\sum_Pm^f_{g,P}=n_g-[g=f]$, $\sum_{g,P}m^f_{g,P}P=(\mu f)'/\mu_2$, for each possible support
$\{g:n_g\ge1\}$ inside an LP-feasible clique. Every system is decided by an exact rational simplex,
with a Farkas vector when it is infeasible. Only $\#3$ and $\#35$ survive, in both cases.
\end{enumerate}

\begin{lemma}[{EQ-LIT claim C5, EQ-AUD A.6}]\label{lem:no3}
Candidate $\#3$ is impossible.
\end{lemma}
The bordering step for one deleted point is \cite[Prop.~5.3, Cor.~5.4]{Gre18} (numbering of
arXiv:1612.03644v3), and \cite[Remark~5.5]{Gre18} records its instances for $75$ and $95$ points. The
rank-$2$ step for two deleted points uses the device of \cite[Lemma~5.24, Thm.~5.25]{GKMS16}.
\begin{proof}
Let $\rho=\kappa=19$ resp.\ $15$ and $c=\rho-5$, so that $\chi_S=(x+5)^{n-d}(x-\rho)^{d-2}(x-c+1)(x-c-1)$.
The matrix $M=-(S+5I)(S-\rho I)$ is positive semidefinite of rank~$2$, with nonzero eigenvalues
$(\rho-1)\cdot6$ and $(\rho+1)\cdot4$ on the $(c-1)$- and $(c+1)$-eigenspaces. Its diagonal is
$5\rho-(n-1)=2$, and for $i\ne j$, $M_{ij}=-(S^2)_{ij}+cS_{ij}$ is even with $|M_{ij}|\le2$. So $M/2$ is a
positive semidefinite $\{0,\pm1\}$ matrix with unit diagonal and rank~$2$. By \cite[Lemma~5.24]{GKMS16}
it equals $s_As_A^{\mathsf T}+s_Bs_B^{\mathsf T}$ for a partition $A\sqcup B$ of the vertices and sign
vectors $s_A,s_B$ supported on $A$, $B$. Comparing eigenvalues gives $|A|=3(\rho-1)$ and
$|B|=2(\rho+1)$, with $s_A$, $s_B$ spanning the $(c-1)$- and $(c+1)$-eigenspaces. After switching,
$S\mathbf 1_A=(c-1)\mathbf 1_A$ and $S\mathbf 1_B=(c+1)\mathbf 1_B$. Bordering $S$ by a vertex with
entries $+1$ on $A$ and $-1$ on $B$ gives a Seidel matrix of order $n+1$: the eigenvectors of $S$ for
$-5$ and $\rho$ are orthogonal to $\mathbf 1_A,\mathbf 1_B$ and stay eigenvectors, and on
$\operatorname{span}\{e_{\mathrm{new}},\mathbf 1_A,\mathbf 1_B\}$ the bordered matrix has characteristic
polynomial $(x-c)(x-\rho)(x+5)$, by $|A|+|B|=n$ and $|A|-|B|=c$. So its characteristic polynomial is
$(x+5)^{n+1-d}(x-c)(x-\rho)^{d-1}$, that is, $(x+5)^{75}(x-14)(x-19)^{19}$ resp.\
$(x+5)^{56}(x-10)(x-15)^{18}$. No Seidel matrix has this characteristic polynomial
\cite[Thm.~5.1]{GSY21}, \cite[Remark~5.5]{Gre18}: by \cite[Cor.~5.4]{Gre18} its switching class would
contain an srg$(95,40,12,20)$ resp.\ srg$(75,32,10,16)$, and these do not exist \cite{AM20,AM18}.
(Equivalently, the bordered matrix would give $n+1$ equiangular lines in $\R^d$, i.e.\ $95$ in $\R^{20}$
resp.\ $75$ in $\R^{19}$, contradicting $N(20)\le94$ \cite[Thm.~5.5]{GSY21} resp.\ $N(19)\le74$
\cite[Thm.~5.2]{GSY21}.)
\end{proof}

So $\chi_S$ is the spectrum $\#35$. We write $\lambda_1>\lambda_2>\lambda_3$ for its eigenvalues other
than $-5$: $(21,19,15)$ with multiplicities $(1,16,3)$ for $94$ lines, and $(17,15,11)$ with
multiplicities $(1,15,3)$ for $74$ lines. Let $m_2=16$ resp.\ $15$ be the multiplicity of $\lambda_2$.

\section{The structure of the spectrum \texorpdfstring{$\#35$}{\#35}}\label{sec:35}

Let $E_\lambda$ be the spectral projections of $S$. Since $G=S+5I=\sum_{\lambda\ne-5}(\lambda+5)E_\lambda$
is the Gram matrix of the lines (up to the factor $5$), each line $v$ has components $E_\lambda e_v$. We
scale them: $X_v=(\sqrt2\,r_v,\sqrt3\,z_v,p_v)\in V_1\oplus V_2\oplus V_3$, where
$V_1\cong\R$, $V_2\cong\R^{m_2}$, $V_3\cong\R^3$ are the eigenspaces of $\lambda_1,\lambda_2,\lambda_3$, and
\[
2r_ur_v=6(\lambda_1+5)(E_1)_{uv},\qquad 3\langle z_u,z_v\rangle=6(\lambda_2+5)(E_2)_{uv},\qquad
\langle p_u,p_v\rangle=6(\lambda_3+5)(E_3)_{uv}.
\]
Then $\langle X_u,X_v\rangle=6G_{uv}$, i.e.\ $2r_ur_v+3\langle z_u,z_v\rangle+\langle p_u,p_v\rangle=6S_{uv}$
for $u\ne v$, and $2r_v^2+3|z_v|^2+|p_v|^2=30$. Since $E_2^2=E_2$ and $E_2$ has rank $m_2$, the matrix
$Z$ with columns $z_v$ satisfies the \emph{frame identity}
\begin{equation}\label{eq:frame}
\sum_vz_vz_v^{\mathsf T}=ZZ^{\mathsf T}=2(\lambda_2+5)\,I_{m_2}\qquad(=40I_{15}\ \text{resp.}\ 48I_{16}),
\end{equation}
and in the same way $\sum_vr_v^2=3(\lambda_1+5)$ and $\sum_vp_vp_v^{\mathsf T}=6(\lambda_3+5)I_3$ (in an
orthonormal basis of $V_3$). Since the $E_\lambda$ are mutually orthogonal, $\sum_vr_vz_v=0$,
$\sum_vr_vp_v=0$ and $\sum_vz_vp_v^{\mathsf T}=0$. Finally $SE_\lambda=\lambda E_\lambda$ gives the
\emph{eigen-equations} $\sum_{u\ne v}S_{vu}(r_u,z_u,p_u)=(\lambda_1r_v,\lambda_2z_v,\lambda_3p_v)$.

\paragraph{Types.} For every vertex $v$, $\chi_{S[v]}=\mu g_v$ with $g_v$ one of the six deck members
$g_0,\dots,g_5$ of $\#35$ that survive level~2, and $(E_\lambda)_{vv}=g_v(\lambda)/\Min'(\lambda)$. This
gives $(r_v^2,|z_v|^2,|p_v|^2)=(0,9,3)$, $(1,8,4)$, $(4,5,7)$, $(5,4,8)$, $(8,1,11)$, $(9,0,12)$ for
types $0,\dots,5$, in both cases. For $u\ne v$, $(E_\lambda)_{uv}=\psi_{uv}(\lambda)/\Min'(\lambda)$ with the
pair polynomial $\psi_{uv}$ of level~2, so every pair of types has a finite list of tuples
$(S_{uv},r_ur_v,\langle z_u,z_v\rangle,\langle p_u,p_v\rangle)$. They are checked mechanically
(\texttt{verify/c35\_structure.py}, checks (1)--(3)). Types $3$ and $4$ have pair polynomials only with
type~$0$ and with themselves.

\paragraph{Supports and counts.} The vertex-star systems of level~2 with the counts $n_t$ of the six
types as unknowns are infeasible except for three supports, and each of these has a unique solution:

\smallskip
\begin{center}
\begin{tabular}{lcc}
\toprule
support & $(n_0,n_1,n_2,n_5)$ for $74$ lines & for $94$ lines\\
\midrule
$(0,1,5)$ & $(24,48,0,2)$ & $(32,60,0,2)$\\
$(0,1,2,5)$ & $(25,45,3,1)$ & $(33,57,3,1)$\\
$(0,1,2)$ & $(26,42,6,0)$ & $(34,54,6,0)$\\
\bottomrule
\end{tabular}
\end{center}
\smallskip
(\texttt{src/star\_counts\_35.py}; EQ-LIT's \texttt{t3\_counts.py} derives the same counts.)

\begin{proposition}[rigidity of $V_1$ and $V_3$]\label{prop:rigid}
After switching, and in a suitable orthonormal frame of $V_3$, every line's $(r_v,p_v)$ depends only on
its \emph{kind}:
\begin{itemize}
\item type $0$ (kind T0): $r=0$, $p=w$; type $1$: $r=1$, $p=2u_x$ for a class $x\in\{a,b,c\}$ (kinds
T1a, T1b, T1c); type $5$: $r=3$, $p=\pm2w$; type $2$: $r=2$, $p=\sigma w-2u_x$ with $\sigma=\pm1$;
\end{itemize}
where $|w|^2=3$, $w\perp u_a,u_b,u_c$, and $u_a,u_b,u_c$ are unit vectors at $120^\circ$
($u_a+u_b+u_c=0$). The numbers $n_x=|L_x|$ of type-$1$ lines of class $x$ and $m_x$ of type-$2$ lines
with $p=\sigma w-2u_x$ are: $(n_x,m_x)=(16,0)$ resp.\ $(20,0)$ in $(0,1,5)$, $(15,1)$ resp.\ $(19,1)$ in
$(0,1,2,5)$, and $(14,2)$ resp.\ $(18,2)$ in $(0,1,2)$, for every $x$. In $(0,1,5)$ the two type-$5$
lines have $p=2w$ and $p=-2w$. In $(0,1,2,5)$ the type-$5$ line has $p=2w$ and every type-$2$ line has
$\sigma=-1$. In $(0,1,2)$ there is one type-$2$ line for each pair $(\sigma,x)$.
\end{proposition}
\begin{proof}
Switching a line negates $X_v$; make $r_v\ge0$ for all $v$, so $r=0,1,2,3$ for types $0,1,2,5$. The
pair data (check (3)) give $\langle p,p'\rangle=\pm3=\pm|p||p'|$ for two type-$0$ lines, so all type-$0$
parts are parallel; switch the type-$0$ lines so that they are all equal to $w$. Type~$0$/type~$5$
gives $\langle p,p'\rangle=\pm6=\pm|p||p'|$, so $p=\pm2w$ for type~$5$. Type~$0$/type~$1$ gives
$\langle p,p'\rangle=0$, so $p=2u$ with $u\perp w$ a unit vector, and type~$1$/type~$1$ gives
$\langle p,p'\rangle\in\{4,-2\}$: the vectors $u$ are equal or at $120^\circ$, so they lie in a set
$\{u_a,u_b,u_c\}$ of three unit vectors at $120^\circ$ in the plane $w^\perp$, $u_a+u_b+u_c=0$.
Type~$0$/type~$2$ gives $\langle p,p'\rangle=\pm3$, so $p=\sigma w+q$ with $q\perp w$, $|q|=2$, and
type~$1$/type~$2$ gives $\langle2u,q\rangle\in\{-4,2\}$ for every type-$1$ vector $u$. The vectors
$q\in w^\perp$ of norm $2$ with $\langle u,q\rangle\in\{-2,1\}$ are exactly $-2u_a,-2u_b,-2u_c$, so
$q=-2u_x$ for one class $x$. On the plane $w^\perp$ the frame
$\sum_vp_vp_v^{\mathsf T}=6(\lambda_3+5)I_3$ reads $\sum_x4(n_x+m_x)u_xu_x^{\mathsf T}=cI_2$. Three unit
vectors at $120^\circ$ form a tight frame only with equal weights (and two or one vectors never do), so
all three classes occur and $n_x+m_x$ is the same for every $x$: $16$ resp.\ $20$. Next
$\sum_vr_vp_v=0$. In $(0,1,5)$ let the two type-$5$ lines have $p=2\varepsilon_1w$ and $p=2\varepsilon_2w$.
Only they contribute to $\langle\sum_vr_vp_v,w\rangle$, which is therefore
$3\cdot2|w|^2(\varepsilon_1+\varepsilon_2)=18(\varepsilon_1+\varepsilon_2)=0$, so $\varepsilon_1=-\varepsilon_2$. In $(0,1,2,5)$, $z_f=0$ for the
type-$5$ line $f$ forces $6S_{tf}=2r_tr_f+\langle p_t,p_f\rangle$ for every line $t$; replacing $w$ by $-w$
(and switching the type-$0$ lines) we may take $p_f=2w$, and then a type-$2$ line has
$S_{tf}=2+\sigma$, so $\sigma=-1$. The component of $\sum_vr_vp_v=0$ in $w^\perp$ reads
$\sum_x(2n_x-4m_x)u_x=0$, so $2n_x-4m_x$ is the same for all $x$; with $n_x+m_x=16$ (resp.\ $20$) and
$\sum m_x=n_2=3$ this gives $m_x=1$. In $(0,1,2)$ the same argument gives $m_x=2$. The $w$-component of
$\sum_vr_vp_v=0$ gives $\sum_t\sigma_t=0$ over the type-$2$ lines, and the $(w,w^\perp)$ block of the
frame gives $\sum_x(\sum_{t\in x}\sigma_t)u_x=0$, so the three class sums of $\sigma$ are equal, hence~$0$.
\end{proof}

\paragraph{The problem in $V_2$.} Put $K(u,v)=2r_ur_v+\langle p_u,p_v\rangle$, a function of the two kinds.
Then
\begin{equation}\label{eq:zz}
3\langle z_u,z_v\rangle=6S_{uv}-K(u,v),
\end{equation}
so the $z_v\in V_2$ have prescribed norms ($9,8,5,0$ for types $0,1,2,5$) and inner products determined
by $S_{uv}=\pm1$, subject to Cauchy--Schwarz. They satisfy the frame identity~\eqref{eq:frame},
$\sum_vr_vz_v=0$ and $\sum_vz_vp_v^{\mathsf T}=0$. The $V_1$- and $V_3$-components of the eigen-equations,
$\sum_{u\ne v}S_{vu}(r_u,p_u)=(\lambda_1r_v,\lambda_3p_v)$, are a $4\times4$ linear system in the sums of the
row $S_{v\cdot}$ over the blocks T0, T1a, T1b, T1c. They fix the number of $-1$ entries of every row in
every block, given the row's entries against the type-$2$ and type-$5$ lines (the \emph{row counts}). The
$V_2$-component gives the block sums $Z_\beta=\sum_{v\in\beta}z_v$ as fixed combinations of the $z$ of
the type-$2$ lines. For example, in $(0,1,5)$ for $74$ lines: two type-$0$ lines have
$\langle z,z'\rangle\in\{1,-3\}$ (per row: $15$ times $1$, $8$ times $-3$); type~$0$/type~$1$:
$\pm2$ ($8$ and $8$ per class); two lines of the same class: $\{0,-4\}$ ($13$ times $0$, twice $-4$); two
type-$1$ lines of different classes: $\pm2$ ($8$ and $8$).

\section{The base-class model}\label{sec:model}

Fix a class $a$; everything is symmetric in $a,b,c$. Let $q_1,\dots,q_{n_a}$ be the $z$-vectors of the
lines of $L_a$, and call a type-$2$ line of class $y\ne a$ an \emph{other-class} type-$2$ line.

\begin{lemma}[the class graph]\label{lem:classgraph}
$\langle q_i,q_j\rangle\in\{0,-4\}$ for $i\ne j$, and $\langle z_t,q_i\rangle\in\{0,-4\}$ for an other-class
type-$2$ line $t$. Every $q_i$ has exactly two $-4$ neighbours among the $q_j$ and the other-class
type-$2$ lines. Hence the $(-4)$-graph on $L_a$ is a disjoint union of paths and cycles, the ends of each
path are joined to other-class type-$2$ lines (a single vertex to two of them), the Gram matrix of $L_a$
is $M=8I-4A$ with $A$ the adjacency matrix of this graph, and all cycles have even length $\ge4$.
\end{lemma}
\begin{proof}
The inner products follow from~\eqref{eq:zz}: $K=2+4=6$ for two class-$a$ lines, and
$K=4+\langle2u_a,\sigma w-2u_y\rangle=6$ for an other-class type-$2$ line. Every other line $y$ with $z_y\ne0$
has $K(y,q_i)=0$, so $\langle z_y,q_i\rangle=2S_{y,q_i}=\pm2$: type $0$ ($K=0$), type $1$ of another class
($K=2-2=0$), and type $2$ of class $a$ ($K=4-4=0$). Now apply~\eqref{eq:frame} to $q_i$:
$\sum_v\langle z_v,q_i\rangle^2=2(\lambda_2+5)\cdot8$. The left side is $64$ (from $v=q_i$), plus $4$ for each
line $y$ of the kinds just listed, plus $16$ for each $-4$ neighbour. With the counts of the table and
of Proposition~\ref{prop:rigid}, the lines of the listed kinds number $n_0+2n_x+m_x$, and
$64+4(n_0+2n_x+m_x)=288$ (for $74$ lines) resp.\ $352$ (for $94$ lines) in all three supports, while
$2(\lambda_2+5)\cdot8=320$ resp.\ $384$. So there are exactly two $-4$ neighbours. A vertex of degree two
in a simple graph lies on a path or a cycle. On a cycle component $C$ the all-ones vector spans the kernel
of $M|_C$, so for every line $y$ the vector $s_i=\langle z_y,q_i\rangle/2$ ($i\in C$) is orthogonal to it;
a type-$0$ line has $s_i=\pm1$, so $|C|$ is even.
\end{proof}

\paragraph{Walks and F-parts.} For a line $y\notin L_a$ write $s(y)=(\langle z_y,q_i\rangle/2)_i$ (the
\emph{steps}; each $s_i$ takes the two values fixed by the kind of $y$ and $S_{y,q_i}$). The projection
$\omega_y$ of $z_y$ onto $\operatorname{span}(L_a)$ satisfies
$\langle\omega_x,\omega_y\rangle=4s(x)^{\mathsf T}M^+s(y)$, and it exists only if $s(y)$ lies in the range of
$M$ (on each cycle: sum zero). The rest $\varphi_y=z_y-\omega_y$ lies in
$F=\operatorname{span}(L_a)^\perp$, of dimension $\dim F=m_2-n_a+(\text{number of cycles})$; a structure with
$\dim F<0$ is impossible. By~\eqref{eq:frame}, $\sum_y\varphi_y\varphi_y^{\mathsf T}=2(\lambda_2+5)I_F$; the
class-$a$ lines have $\varphi=0$ and the type-$5$ lines have $z=0$, so the $\varphi$ of the type-$0$,
type-$1$ (classes $b,c$) and type-$2$ lines span $F$.

\paragraph{References.} Let $D$ be the least common multiple of the denominators of $4M^+$. Since
$(6S-K)/3$ is always an integer, $D\langle\varphi_x,\varphi_y\rangle\in\Z$ for all $x,y$. We choose
\emph{references} $\varphi_{j_1},\dots,\varphi_{j_k}$ as follows. First a maximal independent set of
$F$-parts of type-$2$ lines (their Gram matrix is fixed by the instance). Then repeatedly, while
$k<\dim F$: \emph{case A}, some type-$0$ line has $\varphi$ outside the span of the references; add the
type-$0$ line with the largest residual. \emph{Case B}, otherwise; add the type-$1$ line (class $b$ or
$c$) with the largest residual. Every configuration is reached by this rule (ties included). The later
lines must satisfy the recorded maximality conditions. With $G$ the Gram matrix of the references
(scaled by $D$), every line $y$ gets integer coordinates $C_y=(D\langle\varphi_y,\varphi_{j_r}\rangle)_r$ and
residual $R_y=D\det(G)|\varphi_y|^2-C_y^{\mathsf T}\operatorname{adj}(G)C_y\ge0$, all determined by its kind,
its steps and its entries $S$ against the references. At $k=\dim F$ every $R_y=0$ and
$\langle\varphi_x,\varphi_y\rangle=C_x^{\mathsf T}G^{-1}C_y/D$ exactly. A \emph{node} is a possible
(kind, steps, coordinates) of a line.

\paragraph{Pair rule.} Two nodes $x,y$ can both occur iff
$4s(x)^{\mathsf T}M^+s(y)+\langle\varphi_x,\varphi_y\rangle=(6S-K(x,y))/3$ for $S=+1$ or $S=-1$, and then
$S$ is their Seidel entry. Two distinct lines never have the same node (equal $z$ would need $S=5$
in~\eqref{eq:zz}). So a configuration is a set of pairwise compatible nodes.

\paragraph{Symmetry.} Automorphisms of the class graph that fix the type-$2$ data map configurations to
configurations, and the candidate set of each reference is invariant under the stabiliser of the
references chosen so far. So one candidate per orbit suffices. Every automorphism used is checked to
preserve $M$.

\paragraph{Leaf conditions.} At a leaf (all references chosen) a configuration selects, for each block,
the right number of nodes, and it must satisfy: the frame~\eqref{eq:frame} in walk, walk--$F$ and $F$
coordinates; the block sums; the numbers of $-1$ entries in the columns of the type-$2$ lines; the row
counts of every selected node; and pairwise compatibility. All of these are necessary conditions.

\section{\texorpdfstring{Computation for $\#35$}{Computation for \#35}}\label{sec:comp35}

\paragraph{Program 1} (\texttt{src/c35gen.py}, \texttt{c35k.py}, \texttt{c35t.py}, \texttt{c35u.py};
OR-Tools CP-SAT). It builds explicit coordinates for the class vectors ($2(e_k-e_{k+1})$ along each path or
cycle), computes walks and $F$-coordinates from them, checks the pair rule against explicit vectors
(\texttt{verify/c35\_relation\_test.py}: no mismatch), and decides each leaf with CP-SAT.

\paragraph{Program 2} (\texttt{verify/c35\_cert.py}), written without program~1's code. It takes the
$p$-parts in the basis $(w,u_a,u_b)$, solves the row counts and block sums from the eigen-equations,
uses the class graph only through $M$ (exact Moore--Penrose inverse), enumerates structures, relation
graphs and partner assignments without symmetry reduction, and reduces own type-$2$ lines and
references by explicit automorphisms (each one checked). Every leaf is decided by an LP over $0/1$ node
variables: equality rows for the block sizes, the frame, the block sums, the type-$2$ columns and the row
counts of the selected nodes, with box bounds. Infeasibility is proved by an integer Farkas vector $Y$,
whose margin $\sum_rY_rb_r-\max_{\ell\le x\le u}\sum_rY_ra_rx$ is checked to be positive in exact
arithmetic. Branch and bound over cardinality groups (child $i$ fixes the first $i$ options to $0$ and
option $i$ to $1$; the last child fixes all options to $0$) and propagation (compatibility, a counting
core, cardinality groups) refine the leaf; the propagation depends only on the fixings, and the checker
recomputes it. A certificate file has one record per leaf (instance, reference path, references, leaf
index, tree), and \texttt{--check} re-verifies a whole file without an LP solver, stopping on a missing,
extra or altered record (all three tested). The relation matrix is computed in floating point with
tolerance $10^{-6}$; an exactly compatible pair always passes, so the computed relation contains the true
one, which only weakens the model.

\paragraph{Agreement.} On $984$ identical leaves (all supports, $\dim F\le2$ for $74$ lines) the node sets
and relation matrices of the two programs are equal (\texttt{verify/c35\_cert\_compare.py}). Program~2's
pair rules equal those allowed by the level-$2$ pair data for every pair of kinds
(\texttt{verify/c35\_structure.py}, check~(4)). A third program, written by the team's auditor EQ-AUD
from its own derivation of \S\S\ref{sec:35}--\ref{sec:model} (exact pair relation, a different
reference rule, no shared code), excludes all $444$ leaves of its own search for $94$ lines with
integer Farkas vectors.

\paragraph{Coverage.} In support $(0,1,5)$ there are no type-$2$ lines, so the class graph is a union of
even cycles and $\dim F=m_2-n_a+c$. For $94$ lines ($m_2=16$, $n_a=20$) this needs $c\ge4$ cycles: the
structures are $(8,4,4,4)$, $(6,6,4,4)$ and $(4,4,4,4,4)$. For $74$ lines ($m_2=15$, $n_a=16$) all seven
partitions of $16$ into even parts $\ge4$ occur. In the other supports the structures are the unions of
paths and even cycles allowed by Lemma~\ref{lem:classgraph}.

\begin{center}
\small
\begin{tabular}{llll}
\toprule
case & support & program 2: leaves & program 1\\
\midrule
$94$ & $(0,1,5)$ & $1+1+10$ for $(8,4,4,4)$, $(6,6,4,4)$, $(4,4,4,4,4)$ & all infeasible\\
$94$ & $(0,1,2,5)$ & $3$ (seven other instances: type-$2$ Gram) & all excluded\\
$94$ & $(0,1,2)$ & $14$ (others: type-$2$ Gram, or no own type-$2$ lines) & all excluded\\
$74$ & $(0,1,5)$ & $1+50+28+26+5873+4648$ for the structures & all infeasible\\
& & $(16),(12,4),(10,6),(8,8),(8,4,4),(6,6,4)$ & \\
$74$ & $(0,1,5)$ & $68{,}378+68{,}209$ for $(4,4,4,4)$ & all infeasible (\S\ref{sec:d19})\\
$74$ & $(0,1,2,5)$ & $24{,}792$ & all excluded\\
$74$ & $(0,1,2)$ & $163{,}015$ & all excluded\\
\bottomrule
\end{tabular}
\end{center}

Every program-$2$ leaf in the table is infeasible. All program-$2$ certificate files are re-verified by \texttt{verify/c35\_cert\_checkall.py} with one frozen
version of \texttt{c35\_cert.py} (sha256 prefix \texttt{911efe268df167a3}); the logs are in
\texttt{logs/cert\_final\_check*}.

\section{\texorpdfstring{Proofs of Theorems~\ref{thm:d20} and~\ref{thm:d19}}{Proofs of Theorems 2 and 3}}\label{sec:d19}

\begin{proof}[Proof of Theorem~\ref{thm:d20}]
Suppose there are $94$ equiangular lines in $\R^{20}$. By Lemma~\ref{lem:169}, levels~1 and~2 and
Lemma~\ref{lem:no3}, $\chi_S=(x-21)(x-19)^{16}(x-15)^3(x+5)^{74}$. By \S\ref{sec:35} the lines have one of
the three supports, with the counts of the table and the rigid $(r,p)$-parts of
Proposition~\ref{prop:rigid}. By \S\ref{sec:model} every configuration of the $z$-vectors appears at a
leaf of the base-class model for some structure. Program~2 proves every leaf infeasible with checked
certificates, and program~1 independently finds every leaf infeasible. Every set of more than $94$ lines
contains $94$, so $N(20)\le93$.
\end{proof}

\paragraph{The computations for $74$ lines.} The same chain applies with the $74$-line data. The
largest search is support $(0,1,5)$ with class structure $(4,4,4,4)$, i.e.\ four $4$-cycles, so
$\dim F=3$. For this structure every line $y\notin L_a$ has steps $\pm1$ with sum $0$ on each cycle, and its
projection onto the span of each $4$-cycle is one of six vectors (four of norm $2$, two of norm $1$). Of the
$56$ lines of type $0$ and classes $b,c$, exactly $24$ alternate on each cycle. In particular, if $a$ is
chosen as a class with the fewest cycles, all three classes have structure $(4,4,4,4)$.
The computations (completed 23 October 2026, 13:38 UTC):
\begin{itemize}
\item Supports $(0,1,2,5)$ and $(0,1,2)$, and the six other structures of $(0,1,5)$: excluded by
program~2 (all certificate files checked) and by program~1.
\item Structure $(4,4,4,4)$, program~2: complete. The level-$2$ reference tree has $785$ units in two
shards, with $68{,}378+68{,}209=136{,}587$ leaves, all excluded with certificates. In $42$ of them (units
$217$, $232$ and $240$) the type-$0$ block has hundreds of LP-feasible completions, and each completion
needs a search over the type-$1$ lines of about $1{,}500$ LPs, so the run's budget of $400$ LPs per leaf
was not enough. Branching on the type-$1$ groups first (\texttt{verify/c35\_cert\_hard.py}; certificates
in the same format, so the same check applies to any choice of branching group) closed all $42$, in
$1{,}394$ to $6{,}579$ LPs each. The merged certificate file of each shard passes \texttt{--check} with
the frozen code.
\item Structure $(4,4,4,4)$, program~1: all $1{,}490$ second-level leaves are infeasible in CP-SAT,
$1{,}258$ directly (\texttt{src/c35k.py}, \texttt{src/c35k3.py}) and $232$ after a split by a third
reference ($188{,}648$ sub-leaves, \texttt{src/c35k3.py}). \texttt{verify/c35k3\_coverage.py}
re-enumerates every second-level leaf and its third-level leaves and finds an infeasible record for each
in the union of the run logs; no leaf or sub-leaf was found feasible. With three references every
$F$-part is exact. On $92$ sampled leaves the node sets and pair relations equal program~2's
(\texttt{verify/c35k3\_compare.py}).
\end{itemize}

\begin{proof}[Proof of Theorem~\ref{thm:d19}]
Suppose there are $74$ equiangular lines in $\R^{19}$. By Lemma~\ref{lem:169}, levels~1 and~2 and
Lemma~\ref{lem:no3}, $\chi_S=(x-17)(x-15)^{15}(x-11)^3(x+5)^{55}$. By \S\ref{sec:35} the lines have one of
the three supports, with the counts of the table and the rigid $(r,p)$-parts of
Proposition~\ref{prop:rigid}, and by \S\ref{sec:model} every configuration appears at a leaf of the
base-class model for some structure. By the computations above, program~2 proves every leaf infeasible
with checked certificates, and program~1 independently finds every leaf infeasible. Every set of more
than $74$ lines contains $74$, so $N(19)\le73$.
\end{proof}

\input{t18_sections}

\section{What remains}
\begin{itemize}
\item $344\le N(44)\le420$. For $n=420$, Lemma~\ref{lem:force} forces only $(x-59)^8$ and leaves a
remainder of degree $36$, far beyond direct enumeration.
\item For $d=45,46$ the relative bound ($540$, $736$) is an integer; equality is the existence of
srg$(539,234,81,117)$, resp.\ srg$(735,318,109,159)$. Proposition~\ref{prop:U} does not apply at
$n=R$ and fails at $n=R-1$.
\item $d=20$: $90\le N(20)\le93$. Excluding $93$ lines would need the chain at $n=93$, an odd
order. There Lemma~\ref{lem:force} forces no eigenvalue: the angle is $\arccos\frac15$, $-5$ has
multiplicity $73$, and the best centre $\kappa=19$ gives $T=81$, while $81^{20}>20^{20}4^{19}$. So
the remainder has degree $20$ instead of $13$.
\item $d=19$: $72\le N(19)\le73$. Likewise at $n=73$: $-5$ has multiplicity $54$, the best centre
$\kappa=15$ gives $T=81$, $81^{19}>19^{19}4^{18}$, and the remainder has degree $19$. (For program~2's
$42$ hard leaves of \S\ref{sec:d19}, lookahead branching, conditional row-count cuts, eigen-equation
rows for selected nodes, $4$-cycle propagation, spectral cuts and CP-SAT did not help within the
default branching order; branching on the type-$1$ groups first decided all of them.)
\item $d=18$: $57\le N(18)\le58$. At $n=58$ (even order), $-5$ has multiplicity $40$, the best centre
$\kappa=11$ gives $T=84>4\cdot18$, so Lemma~\ref{lem:force} forces no eigenvalue, and the remainder
has degree $18$ instead of $17$.
\item \emph{Size of the next enumerations} (estimates only: systematic samples of about one
subtree in $1{,}000$ at depth $5$ of the enumeration tree of Lemma~\ref{lem:59enum}, run with the
same congruences and filters but without a forced eigenvalue; scratch code \texttt{src/est/},
samples in \texttt{runs/est\_enum/}). On $59$ lines itself the sample predicts about $1{,}000$
candidates, against the true $1{,}214$. For $73$ lines in $\R^{19}$ and $93$ in $\R^{20}$ the
samples suggest $2\cdot10^9$ and $1.5\cdot10^9$ nodes, about $1{,}300$ and $2{,}300$
candidates, and $170$ and $130$ CPU-hours for our enumerator, against $4\cdot10^7$ nodes and $16$
CPU-hours for $59$ lines. Every sampled candidate has the centre as an eigenvalue ($15$ with
multiplicity $9$ to $12$, resp.\ $19$ with multiplicity $3$ to $13$). For $58$ lines the tree is
far larger, of order $10^{10}$ nodes with $10^5$ to $10^6$ leaves: at even order the classes
modulo $2^7$ are exactly the type-$2$ residues (a sample with and without the filter of
$P_{58,7}$ visits the same nodes), and \cite[Lemma~3.12]{GY19} needs odd order, so only
real-rootedness prunes. Every sampled leaf has the eigenvalue $11$.
\item Dimensions $18$ and $42$: the moment tests give nothing new.
\end{itemize}

\begin{thebibliography}{99}
\bibitem{AM18} J.~Azarija, T.~Marc, There is no $(75,32,10,16)$ strongly regular graph, Linear Algebra Appl.\ 557 (2018) 62--83.
\bibitem{AM20} J.~Azarija, T.~Marc, There is no $(95,40,12,20)$ strongly regular graph, J.\ Combin.\ Des.\ 28 (2020) 294--306.
\bibitem{BY14} A.~Barg, W.-H.~Yu, New bounds for equiangular lines, in: Discrete Geometry and Algebraic Combinatorics, Contemp.\ Math.\ 625, AMS (2014) 111--121.
\bibitem{CKLY22} M.-Y.~Cao, J.~H.~Koolen, Y.-C.~R.~Lin, W.-H.~Yu, The Lemmens--Seidel conjecture and forbidden subgraphs, J.\ Combin.\ Theory Ser.\ A 185 (2022) 105538.
\bibitem{GKMS16} G.~Greaves, J.~H.~Koolen, A.~Munemasa, F.~Sz\"oll\H{o}si, Equiangular lines in Euclidean spaces, J.\ Combin.\ Theory Ser.\ A 138 (2016) 208--235.
\bibitem{Gre18} G.~Greaves, Equiangular line systems and switching classes containing regular graphs, Linear Algebra Appl.\ 536 (2018) 31--51.
\bibitem{GS24} G.~R.~W.~Greaves, J.~Syatriadi, Real equiangular lines in dimension $18$ and the Jacobi identity for complementary subgraphs, J.\ Combin.\ Theory Ser.\ A 201 (2024) 105812.
\bibitem{GS25} G.~R.~W.~Greaves, J.~Syatriadi, Real-rooted integer polynomial enumeration algorithms and interlacing polynomials via linear programming, arXiv:2504.09241 (2025).
\bibitem{GSY21} G.~R.~W.~Greaves, J.~Syatriadi, P.~Yatsyna, Equiangular lines in low dimensional Euclidean spaces, Combinatorica 41 (2021) 839--872. We cite the numbering of arXiv:2002.08085v1.
\bibitem{GSY23} G.~R.~W.~Greaves, J.~Syatriadi, P.~Yatsyna, Equiangular lines in Euclidean spaces: dimensions 17 and 18, Math.\ Comp.\ 92 (2023) 1867--1903.
\bibitem{GY19} G.~R.~W.~Greaves, P.~Yatsyna, On equiangular lines in 17 dimensions and the characteristic polynomial of a Seidel matrix, Math.\ Comp.\ 88 (2019) 3041--3061.
\bibitem{KT19} E.~J.~King, X.~Tang, New upper bounds for equiangular lines by pillar decomposition, SIAM J.\ Discrete Math.\ 33 (2019) 2479--2508.
\bibitem{LS73} P.~W.~H.~Lemmens, J.~J.~Seidel, Equiangular lines, J.\ Algebra 24 (1973) 494--512.
\bibitem{Szo19} F.~Sz\"oll\H{o}si, A remark on a construction of D.~S.~Asche, Discrete Comput.\ Geom.\ 61 (2019) 120--122.
\bibitem{Tay71} D.~E.~Taylor, Some topics in the theory of finite groups, PhD thesis, University of Oxford (1971).
\bibitem{Taylor} D.~E.~Taylor, Regular 2-graphs, Proc.\ London Math.\ Soc.\ (3) 35 (1977) 257--274.
\bibitem{vLS66} J.~H.~van Lint, J.~J.~Seidel, Equilateral point sets in elliptic geometry, Indag.\ Math.\ 28 (1966) 335--348.
\bibitem{Yu17} W.-H.~Yu, New bounds for equiangular lines and spherical two-distance sets, SIAM J.\ Discrete Math.\ 31 (2017) 908--917.
\end{thebibliography}

\appendix
\section{\texorpdfstring{The 82 candidates for 421 lines in $\R^{44}$}{The 82 candidates for 421 lines in R\textasciicircum 44}}\label{app:cand}
Each row gives: the eigenvalues of $\chi_S$ other than $-7^{377}$ (integer eigenvalues with
multiplicity, then irreducible non-linear factors), $\delta$, the deck size, and the number of decimal
digits of the largest certificate entry. The full data are in
\texttt{data/n421/decks/$\langle$idx$\rangle$.json}.
{\small
\input{tables/n421_candidates.tex}
}

\section{\texorpdfstring{Excluded top values for $a\le13$}{Excluded top values for a <= 13}}\label{app:U}
Columns: $a$, $d$, $R$, $n_0=\lfloor R\rfloor$, the smallest odd witness $\kappa$, $s$, $T$, $\eta^*(T)$. In
each row, $N_{1/a}(d)\le n_0-1$. Produced by \texttt{verify/c1\_uniform.py} and
\texttt{verify/c1\_uniform.gp}, whose outputs agree line by line.
{\small
\input{tables/thmU_rows.tex}
}
\end{document}
